\section{Conclusion}

This work introduces the powerful $\DiDI$-technique and solves three open problems in PIT for depth-$4$ circuits, namely $\SPSP{[k]}{}{[\delta]}$ (blackbox) and $\SPSU{[k]}{}{}$ (both whitebox and blackbox). Here are some immediate questions of interest which require rigorous investigation.
\begin{enumerate}
	\item Can the exponent in \autoref{thm:thm1} be improved to $O(k)$? Currently, it is exponential in $k$.    
    \item Can we improve \autoref{thm:thm2}(b) to $s^{O(\log \log s)}$ (like in \autoref{thm:thm2}(a))?
    \item Can we design a polynomial-time PIT for $\SPSP{[k]}{}{[\delta]}$?
    \item Design a polynomial time PIT for $\SWSP{}{}{[\delta]}$ circuits (i.e.~unbounded top-fanin)? 
    \item Can we solve PIT for $\Sigma^{[k]}\Pi\Sigma\mathsf{M}_2$ circuits efficiently (polynomial/quasipolynomial-time), where \(\Sigma\mathsf{M}_2\) denotes bivariate polynomials?    
    \item Can we design an efficient PIT for rational functions of the form $\Sigma^{}\,(1/\Sigma\wedge\Sigma)$ or $\Sigma^{}\,(1/\Sigma\Pi)$ (for {\em un}bounded top-fanin)?
\end{enumerate}
